\section{Рабочий проект}

В процессе разработки были реализованы три самостоятельных Django-приложения модуля взаимодействия с пользователем: \texttt{chats}, \texttt{notifications} и \texttt{contracts}. Каждое приложение включает модели данных, сериализаторы, представления REST API, а также WebSocket-компоненты или сервисные функции в зависимости от назначения.

\subsection{Приложение chats}

\texttt{chats} --- приложение, реализующее подсистему обмена сообщениями между пользователями платформы. Поддерживает создание диалогов с привязкой к объявлению, отправку текстовых сообщений и изображений, доставку новых сообщений в режиме реального времени через WebSocket, а также отметку сообщений как прочитанных. Приложение состоит из моделей данных (\texttt{models.py}), представлений REST API (\texttt{views.py}), WebSocket-консьюмера (\texttt{consumers.py}), сериализаторов (\texttt{serializers.py}) и вспомогательного модуля (\texttt{utils.py}).

Описание моделей данных приложения \texttt{chats} представлено в таблице~\ref{tab:chats_models}.

\begin{xltabular}{\textwidth}{|>{\small\raggedright\arraybackslash}p{0.30\textwidth}|>{\small\raggedright\arraybackslash}p{0.22\textwidth}|>{\small\raggedright\arraybackslash}X|}
\caption{Описание моделей данных приложения chats \label{tab:chats_models}}\\ \hline
\centrow Модель / поле & \centrow Тип & \centrow Описание \\ \hline
\endfirsthead
\continuecaption{Продолжение таблицы \ref{tab:chats_models}}
\centrow Модель / поле & \centrow Тип & \centrow Описание \\ \hline
\finishhead
\leftrow Chat & \leftrow Модель & \leftrow Диалог между двумя пользователями платформы. \\ \hline
\leftrow users & \leftrow ManyToManyField & \leftrow Участники диалога. \\ \hline
\leftrow item & \leftrow ForeignKey (Item) & \leftrow Объявление, к которому привязан чат. Необязательное поле. \\ \hline
\leftrow created\_at & \leftrow DateTimeField & \leftrow Дата и время создания диалога. \\ \hline
\leftrow Message & \leftrow Модель & \leftrow Сообщение, отправленное в рамках диалога. \\ \hline
\leftrow chat & \leftrow ForeignKey (Chat) & \leftrow Диалог, к которому относится сообщение. \\ \hline
\leftrow sender & \leftrow ForeignKey (User) & \leftrow Пользователь, отправивший сообщение. \\ \hline
\leftrow text & \leftrow TextField & \leftrow Текстовое содержимое сообщения. \\ \hline
\leftrow image & \leftrow ImageField & \leftrow Прикреплённое изображение. Необязательное поле. \\ \hline
\leftrow created\_at & \leftrow DateTimeField & \leftrow Дата и время отправки сообщения. \\ \hline
\leftrow is\_read & \leftrow BooleanField & \leftrow Признак прочтения сообщения получателем. \\ \hline
\end{xltabular}

Описание вспомогательных функций модуля \texttt{views.py} приложения \texttt{chats} представлено в таблице~\ref{tab:chats_helpers}.

\begin{xltabular}{\textwidth}{|>{\small\raggedright\arraybackslash}p{0.30\textwidth}|>{\small\raggedright\arraybackslash}p{0.22\textwidth}|>{\small\raggedright\arraybackslash}X|}
\caption{Описание вспомогательных функций views.py приложения chats \label{tab:chats_helpers}}\\ \hline
\centrow Название функции & \centrow Аргументы & \centrow Описание \\ \hline
\endfirsthead
\continuecaption{Продолжение таблицы \ref{tab:chats_helpers}}
\centrow Название функции & \centrow Аргументы & \centrow Описание \\ \hline
\finishhead
\leftrow mark\_chat\_messages\_read & \leftrow chat, user & \leftrow Находит непрочитанные сообщения в чате, не отправленные текущим пользователем, помечает их как прочитанные и возвращает список их идентификаторов. \\ \hline
\leftrow broadcast\_messages\_read & \leftrow chat\_id, message\_ids, read\_by\_user\_id & \leftrow Рассылает через channel layer группе чата событие \texttt{messages\_read} с идентификаторами прочитанных сообщений и идентификатором пользователя. \\ \hline
\end{xltabular}

Описание представлений REST API приложения \texttt{chats} представлено в таблице~\ref{tab:chats_views}.

\begin{xltabular}{\textwidth}{|>{\small\raggedright\arraybackslash}p{0.30\textwidth}|>{\small\raggedright\arraybackslash}p{0.22\textwidth}|>{\small\raggedright\arraybackslash}X|}
\caption{Описание представлений REST API приложения chats \label{tab:chats_views}}\\ \hline
\centrow Класс / метод & \centrow HTTP-метод, маршрут & \centrow Описание \\ \hline
\endfirsthead
\continuecaption{Продолжение таблицы \ref{tab:chats_views}}
\centrow Класс / метод & \centrow HTTP-метод, маршрут & \centrow Описание \\ \hline
\finishhead
\leftrow ConversationListCreateView.\allowbreak get & \leftrow GET /conversations/ & \leftrow Возвращает список чатов текущего пользователя, аннотированных датой последнего сообщения, упорядоченных по убыванию активности. \\ \hline
\leftrow ConversationListCreateView.\allowbreak post & \leftrow POST /conversations/ & \leftrow Создаёт новый диалог с указанным пользователем и опциональной привязкой к объявлению. Если чат уже существует, возвращает его. \\ \hline
\leftrow ConversationDetailView.\allowbreak get & \leftrow GET /conversations/\{id\}/ & \leftrow Возвращает содержимое чата со списком сообщений. При открытии помечает входящие сообщения как прочитанные и рассылает событие прочтения. \\ \hline
\leftrow ConversationMessageCreateView.\allowbreak post & \leftrow POST /conversations/\{id\}/\allowbreak messages/ & \leftrow Сохраняет новое сообщение (текст и/или изображение), создаёт уведомление получателю и рассылает событие \texttt{chat\_message} через WebSocket. \\ \hline
\leftrow MarkReadView.\allowbreak post & \leftrow POST /conversations/\{id\}/\allowbreak mark\_read/ & \leftrow Явно помечает непрочитанные сообщения чата как прочитанные и рассылает событие через WebSocket. \\ \hline
\leftrow SupportFAQView.\allowbreak get & \leftrow GET /support-faq/ & \leftrow Возвращает структурированный набор часто задаваемых вопросов с пошаговыми инструкциями. \\ \hline
\end{xltabular}

Описание сериализаторов приложения \texttt{chats} представлено в таблице~\ref{tab:chats_serializers}.

\begin{xltabular}{\textwidth}{|>{\small\raggedright\arraybackslash}p{0.30\textwidth}|>{\small\raggedright\arraybackslash}p{0.22\textwidth}|>{\small\raggedright\arraybackslash}X|}
\caption{Описание сериализаторов приложения chats \label{tab:chats_serializers}}\\ \hline
\centrow Сериализатор / поле & \centrow Тип & \centrow Описание \\ \hline
\endfirsthead
\continuecaption{Продолжение таблицы \ref{tab:chats_serializers}}
\centrow Сериализатор / поле & \centrow Тип & \centrow Описание \\ \hline
\finishhead
\leftrow MessageSerializer & \leftrow ModelSerializer & \leftrow Сериализатор для чтения сообщений. Включает идентификатор и имя отправителя, текст, ссылку на изображение, дату и признак прочтения. \\ \hline
\leftrow MessageCreateSerializer & \leftrow ModelSerializer & \leftrow Сериализатор для создания сообщения. Проверяет, что хотя бы одно из полей \texttt{text} или \texttt{image} заполнено. \\ \hline
\leftrow ChatSerializer & \leftrow ModelSerializer & \leftrow Сериализатор чата. Вычисляемые поля: \texttt{other\_user} (данные собеседника), \texttt{unread\_count} (число непрочитанных), \texttt{last\_message} (последнее сообщение), \texttt{item} (связанное объявление). \\ \hline
\end{xltabular}

Описание методов WebSocket-консьюмера \texttt{ChatConsumer} приложения \texttt{chats} представлено в таблице~\ref{tab:chats_consumer}.

\begin{xltabular}{\textwidth}{|>{\small\raggedright\arraybackslash}p{0.30\textwidth}|>{\small\raggedright\arraybackslash}p{0.22\textwidth}|>{\small\raggedright\arraybackslash}X|}
\caption{Описание методов ChatConsumer приложения chats \label{tab:chats_consumer}}\\ \hline
\centrow Метод & \centrow Аргументы & \centrow Описание \\ \hline
\endfirsthead
\continuecaption{Продолжение таблицы \ref{tab:chats_consumer}}
\centrow Метод & \centrow Аргументы & \centrow Описание \\ \hline
\finishhead
\leftrow connect & \leftrow --- & \leftrow Устанавливает WebSocket-соединение. Проверяет аутентификацию пользователя, существование чата и членство в нём. Добавляет канал в группу чата. \\ \hline
\leftrow disconnect & \leftrow close\_code & \leftrow Удаляет канал из группы чата при разрыве соединения. \\ \hline
\leftrow receive & \leftrow text\_data & \leftrow Обрабатывает входящее сообщение от клиента. При действии \texttt{mark\_read} помечает сообщения прочитанными. При наличии текста сохраняет сообщение и рассылает его группе. \\ \hline
\leftrow chat\_message & \leftrow event & \leftrow Доставляет клиенту событие \texttt{new\_message} с данными нового сообщения. \\ \hline
\leftrow messages\_read & \leftrow event & \leftrow Доставляет клиенту событие \texttt{messages\_read} с идентификаторами прочитанных сообщений. \\ \hline
\leftrow save\_message & \leftrow user, text & \leftrow Сохраняет текстовое сообщение в базу данных и инициирует отправку уведомления получателю. \\ \hline
\leftrow serialize\_message & \leftrow message & \leftrow Преобразует объект сообщения в словарь для передачи через WebSocket. \\ \hline
\leftrow chat\_exists & \leftrow --- & \leftrow Проверяет существование чата по идентификатору из URL. \\ \hline
\leftrow is\_chat\_member & \leftrow --- & \leftrow Проверяет, является ли текущий пользователь участником чата. \\ \hline
\leftrow mark\_messages\_read & \leftrow user & \leftrow Помечает непрочитанные входящие сообщения как прочитанные и возвращает их идентификаторы. \\ \hline
\end{xltabular}

\subsection{Приложение notifications}

\texttt{notifications} --- приложение, реализующее двухканальную систему уведомлений пользователей. Первый канал --- всплывающие уведомления на сайте, доставляемые через персональное WebSocket-соединение в реальном времени. Второй канал --- email-оповещения, отправляемые через SMTP при ключевых событиях платформы. Приложение состоит из модели данных (\texttt{models.py}), сервисного слоя (\texttt{services.py}), функций отправки email (\texttt{email.py}), представлений REST API (\texttt{views.py}), WebSocket-консьюмера (\texttt{consumers.py}), сериализатора (\texttt{serializers.py}) и управляющей команды (\texttt{send\_return\_reminders}).

Описание модели данных приложения \texttt{notifications} представлено в таблице~\ref{tab:notif_model}.

\begin{xltabular}{\textwidth}{|>{\small\raggedright\arraybackslash}p{0.30\textwidth}|>{\small\raggedright\arraybackslash}p{0.22\textwidth}|>{\small\raggedright\arraybackslash}X|}
\caption{Описание модели данных приложения notifications \label{tab:notif_model}}\\ \hline
\centrow Модель / поле & \centrow Тип & \centrow Описание \\ \hline
\endfirsthead
\continuecaption{Продолжение таблицы \ref{tab:notif_model}}
\centrow Модель / поле & \centrow Тип & \centrow Описание \\ \hline
\finishhead
\leftrow Notification & \leftrow Модель & \leftrow Уведомление, адресованное конкретному пользователю платформы. \\ \hline
\leftrow \quad user & \leftrow ForeignKey (User) & \leftrow Получатель уведомления. \\ \hline
\leftrow \quad type & \leftrow CharField & \leftrow Тип события. Допустимые значения: \texttt{booking\_created}, \texttt{booking\_confirmed}, \texttt{booking\_cancelled}, \texttt{payment\_confirmed}, \texttt{return\_reminder}, \texttt{new\_message}, \texttt{new\_review}, \texttt{item\_moderation\_request}, \texttt{item\_moderation\_approved}, \texttt{item\_moderation\_rejected}. \\ \hline
\leftrow \quad message & \leftrow TextField & \leftrow Текст уведомления, отображаемый пользователю. \\ \hline
\leftrow \quad metadata & \leftrow JSONField & \leftrow Дополнительные данные для навигации на клиенте: направление перехода, идентификаторы бронирования, объявления, чата. \\ \hline
\leftrow \quad is\_read & \leftrow BooleanField & \leftrow Признак прочтения уведомления. \\ \hline
\leftrow \quad created\_at & \leftrow DateTimeField & \leftrow Дата и время создания уведомления. \\ \hline
\end{xltabular}

Описание функций сервисного модуля \texttt{services.py} приложения \texttt{notifications} представлено в таблице~\ref{tab:notif_services}.

\begin{xltabular}{\textwidth}{|>{\small\raggedright\arraybackslash}p{0.30\textwidth}|>{\small\raggedright\arraybackslash}p{0.22\textwidth}|>{\small\raggedright\arraybackslash}X|}
\caption{Описание функций модуля services.py приложения notifications \label{tab:notif_services}}\\ \hline
\centrow Название функции & \centrow Аргументы & \centrow Описание \\ \hline
\endfirsthead
\continuecaption{Продолжение таблицы \ref{tab:notif_services}}
\centrow Название функции & \centrow Аргументы & \centrow Описание \\ \hline
\finishhead
\leftrow create\_notification & \leftrow user, notification\_type, message, metadata & \leftrow Создаёт запись уведомления в базе данных и немедленно доставляет его пользователю через WebSocket, вызывая \texttt{broadcast\_notification}. \\ \hline
\leftrow broadcast\_notification & \leftrow notification & \leftrow Отправляет сериализованный объект уведомления в персональную channel-группу пользователя. Ошибки доставки не прерывают сохранение уведомления. \\ \hline
\leftrow notify\_booking\_created & \leftrow booking & \leftrow Создаёт уведомление арендодателю о новой заявке на бронирование и отправляет соответствующее email-сообщение. \\ \hline
\leftrow notify\_booking\_confirmed & \leftrow booking & \leftrow Создаёт уведомление арендатору о подтверждении бронирования и отправляет соответствующее email-сообщение. \\ \hline
\leftrow notify\_payment\_confirmed & \leftrow booking & \leftrow Создаёт уведомление арендодателю о поступлении оплаты и отправляет соответствующее email-сообщение. \\ \hline
\leftrow notify\_booking\_cancelled & \leftrow booking, cancelled\_by & \leftrow Создаёт уведомления обеим сторонам сделки об отмене бронирования с указанием инициатора отмены и отправляет email-оповещения. \\ \hline
\leftrow notify\_return\_reminder & \leftrow booking & \leftrow Создаёт уведомление арендатору с напоминанием о необходимости возврата вещи и отправляет соответствующее email-сообщение. \\ \hline
\leftrow notify\_new\_message & \leftrow message & \leftrow Создаёт уведомления и отправляет email-оповещения всем участникам чата, кроме отправителя сообщения. \\ \hline
\leftrow notify\_new\_review & \leftrow review & \leftrow Создаёт уведомление арендодателю о новом отзыве с оценкой и комментарием, а также отправляет email-оповещение. \\ \hline
\end{xltabular}

Описание функций модуля \texttt{email.py} приложения \texttt{notifications} представлено в таблице~\ref{tab:notif_email}.

\begin{xltabular}{\textwidth}{|>{\small\raggedright\arraybackslash}p{0.30\textwidth}|>{\small\raggedright\arraybackslash}p{0.22\textwidth}|>{\small\raggedright\arraybackslash}X|}
\caption{Описание функций модуля email.py приложения notifications \label{tab:notif_email}}\\ \hline
\centrow Название функции & \centrow Аргументы & \centrow Описание \\ \hline
\endfirsthead
\continuecaption{Продолжение таблицы \ref{tab:notif_email}}
\centrow Название функции & \centrow Аргументы & \centrow Описание \\ \hline
\finishhead
\leftrow \_send & \leftrow to\_email, subject, message & \leftrow Базовая функция отправки email через Django \texttt{send\_mail}. При пустом адресе или недоступности почтового сервера молча завершается, не прерывая основной процесс. \\ \hline
\leftrow send\_password\_\allowbreak reset\_code & \leftrow user, code & \leftrow Отправляет пользователю шестизначный код для сброса пароля с указанием срока действия кода. \\ \hline
\leftrow send\_booking\_created & \leftrow booking & \leftrow Отправляет арендодателю email о новой заявке на бронирование с указанием арендатора, дат и суммы аренды. \\ \hline
\leftrow send\_booking\_confirmed & \leftrow booking & \leftrow Отправляет арендатору email о подтверждении бронирования с датами и итоговой суммой. \\ \hline
\leftrow send\_payment\_confirmed & \leftrow booking & \leftrow Отправляет арендодателю email о поступлении оплаты от арендатора с деталями бронирования. \\ \hline
\leftrow send\_booking\_cancelled & \leftrow booking, cancelled\_by & \leftrow Отправляет обеим сторонам email об отмене бронирования с указанием инициатора. \\ \hline
\leftrow send\_return\_reminder & \leftrow booking & \leftrow Отправляет арендатору email с напоминанием о необходимости возврата вещи на следующий день. \\ \hline
\leftrow send\_new\_message & \leftrow recipient, sender\_username & \leftrow Отправляет получателю email с уведомлением о новом сообщении в чате. \\ \hline
\leftrow send\_new\_review & \leftrow booking & \leftrow Отправляет арендодателю email с уведомлением о новом отзыве о товаре. \\ \hline
\leftrow send\_item\_moderation\_\allowbreak request\_to\_\allowbreak admins & \leftrow moderation\_\allowbreak request, admins & \leftrow Отправляет всем администраторам email с деталями новой заявки на модерацию объявления. \\ \hline
\leftrow send\_item\_moderation\_\allowbreak approved & \leftrow user, item\_title & \leftrow Отправляет пользователю email об одобрении и публикации объявления. \\ \hline
\leftrow send\_item\_moderation\_\allowbreak rejected & \leftrow user, item\_title, reason & \leftrow Отправляет пользователю email об отклонении заявки на публикацию объявления с указанием причины. \\ \hline
\end{xltabular}

Описание представлений REST API приложения \texttt{notifications} представлено в таблице~\ref{tab:notif_views}.

\begin{xltabular}{\textwidth}{|>{\small\raggedright\arraybackslash}p{0.30\textwidth}|>{\small\raggedright\arraybackslash}p{0.22\textwidth}|>{\small\raggedright\arraybackslash}X|}
\caption{Описание представлений REST API приложения notifications \label{tab:notif_views}}\\ \hline
\centrow Метод & \centrow HTTP-метод, маршрут & \centrow Описание \\ \hline
\endfirsthead
\continuecaption{Продолжение таблицы \ref{tab:notif_views}}
\centrow Метод & \centrow HTTP-метод, маршрут & \centrow Описание \\ \hline
\finishhead
\leftrow NotificationViewSet.list & \leftrow GET /notifications/ & \leftrow Возвращает список уведомлений текущего пользователя в порядке убывания даты создания. \\ \hline
\leftrow NotificationViewSet.retrieve & \leftrow GET /notifications/\{id\}/ & \leftrow Возвращает данные одного уведомления по идентификатору. \\ \hline
\leftrow NotificationViewSet.\allowbreak unread\_count & \leftrow GET /notifications/\allowbreak unread\_count/ & \leftrow Возвращает количество непрочитанных уведомлений текущего пользователя. \\ \hline
\leftrow NotificationViewSet.\allowbreak mark\_read & \leftrow POST /notifications/\{id\}/\allowbreak mark\_read/ & \leftrow Помечает указанное уведомление как прочитанное и возвращает его обновлённые данные. \\ \hline
\leftrow NotificationViewSet.\allowbreak mark\_all\_read & \leftrow POST /notifications/\allowbreak mark\_all\_read/ & \leftrow Помечает все непрочитанные уведомления пользователя как прочитанные и возвращает количество обновлённых записей. \\ \hline
\end{xltabular}

Описание методов WebSocket-консьюмера \texttt{NotificationConsumer} приложения \texttt{notifications} представлено в таблице~\ref{tab:notif_consumer}.

\begin{xltabular}{\textwidth}{|>{\small\raggedright\arraybackslash}p{0.30\textwidth}|>{\small\raggedright\arraybackslash}p{0.22\textwidth}|>{\small\raggedright\arraybackslash}X|}
\caption{Описание методов NotificationConsumer приложения notifications \label{tab:notif_consumer}}\\ \hline
\centrow Метод & \centrow Аргументы & \centrow Описание \\ \hline
\endfirsthead
\continuecaption{Продолжение таблицы \ref{tab:notif_consumer}}
\centrow Метод & \centrow Аргументы & \centrow Описание \\ \hline
\finishhead
\leftrow connect & \leftrow --- & \leftrow Устанавливает WebSocket-соединение. Проверяет аутентификацию пользователя. Добавляет канал в персональную группу \texttt{notifications\_\{user\_id\}}. \\ \hline
\leftrow disconnect & \leftrow close\_code & \leftrow Удаляет канал из персональной группы при разрыве соединения. \\ \hline
\leftrow notification\_created & \leftrow event & \leftrow Доставляет клиенту событие \texttt{notification\_created} с полными данными нового уведомления. \\ \hline
\end{xltabular}

Управляющая команда \texttt{send\_return\_reminders} реализована в виде Django management command и предназначена для планировщика задач. При запуске выбирает из базы данных все подтверждённые бронирования с датой окончания, приходящейся на следующий день, и для каждого вызывает \texttt{notify\_return\_reminder}, отправляя email и создавая сайт-уведомление арендатору.

\subsection{Приложение contracts}

\texttt{contracts} --- приложение, реализующее подсистему автоматической генерации договоров аренды и их демонстрационного подписания. При первом обращении к договору по подтверждённому бронированию система автоматически формирует уникальный номер документа и создаёт PDF-файл, содержащий полный текст договора с реквизитами сторон. После подписания договора обеими сторонами PDF перегенерируется с заполненными блоками электронной подписи. Приложение состоит из модели данных с методами генерации PDF (\texttt{models.py}), представлений REST API (\texttt{views.py}) и сериализаторов (\texttt{serializers.py}).

Описание вспомогательных функций модуля \texttt{models.py} приложения \texttt{contracts} представлено в таблице~\ref{tab:contracts_helpers}.

\begin{xltabular}{\textwidth}{|>{\small\raggedright\arraybackslash}p{0.30\textwidth}|>{\small\raggedright\arraybackslash}p{0.22\textwidth}|>{\small\raggedright\arraybackslash}X|}
\caption{Описание вспомогательных функций models.py приложения contracts \label{tab:contracts_helpers}}\\ \hline
\centrow Название функции & \centrow Аргументы & \centrow Описание \\ \hline
\endfirsthead
\continuecaption{Продолжение таблицы \ref{tab:contracts_helpers}}
\centrow Название функции & \centrow Аргументы & \centrow Описание \\ \hline
\finishhead
\leftrow \_ensure\_cyrillic\_fonts & \leftrow --- & \leftrow Последовательно проверяет наличие кандидатов кириллических TTF-шрифтов на файловой системе. Регистрирует первый доступный шрифт в ReportLab и возвращает имена обычного, жирного и курсивного начертаний. \\ \hline
\leftrow \_party\_description & \leftrow user & \leftrow Формирует строку описания стороны договора и словарь реквизитов в зависимости от организационно-правовой формы пользователя: физическое лицо, индивидуальный предприниматель или юридическое лицо. \\ \hline
\leftrow \_eds\_block & \leftrow signer\_name, signed\_at, signature\_code, role\_label & \leftrow Формирует строку блока электронной цифровой подписи с именем подписанта, датой и кодом подписи либо статус «Не подписан» при отсутствии данных. \\ \hline
\end{xltabular}

Описание полей и методов модели \texttt{Contract} приложения \texttt{contracts} представлено в таблице~\ref{tab:contracts_model}.

\begin{xltabular}{\textwidth}{|>{\small\raggedright\arraybackslash}p{0.30\textwidth}|>{\small\raggedright\arraybackslash}p{0.22\textwidth}|>{\small\raggedright\arraybackslash}X|}
\caption{Описание полей и методов модели Contract \label{tab:contracts_model}}\\ \hline
\centrow Поле / метод & \centrow Тип / аргументы & \centrow Описание \\ \hline
\endfirsthead
\continuecaption{Продолжение таблицы \ref{tab:contracts_model}}
\centrow Поле / метод & \centrow Тип / аргументы & \centrow Описание \\ \hline
\finishhead
\leftrow booking & \leftrow OneToOneField (Booking) & \leftrow Бронирование, к которому относится договор. \\ \hline
\leftrow document\_number & \leftrow CharField & \leftrow Уникальный номер договора в формате \texttt{AR-\{год\}-\{booking\_id\}-\{UUID\}}. Генерируется автоматически при первом сохранении. \\ \hline
\leftrow file & \leftrow FileField & \leftrow PDF-файл договора, хранящийся в директории \texttt{contracts/}. \\ \hline
\leftrow is\_signed & \leftrow BooleanField & \leftrow Признак полного подписания договора обеими сторонами. \\ \hline
\leftrow renter\_signer\_name & \leftrow CharField & \leftrow Имя подписанта со стороны арендатора. \\ \hline
\leftrow renter\_signature\_code & \leftrow CharField & \leftrow Код демонстрационной электронной подписи арендатора. \\ \hline
\leftrow renter\_signed\_at & \leftrow DateTimeField & \leftrow Дата и время подписания договора арендатором. \\ \hline
\leftrow owner\_signer\_name & \leftrow CharField & \leftrow Имя подписанта со стороны арендодателя. \\ \hline
\leftrow owner\_signature\_code & \leftrow CharField & \leftrow Код демонстрационной электронной подписи арендодателя. \\ \hline
\leftrow owner\_signed\_at & \leftrow DateTimeField & \leftrow Дата и время подписания договора арендодателем. \\ \hline
\leftrow signed\_at & \leftrow DateTimeField & \leftrow Дата и время полного подписания договора обеими сторонами. \\ \hline
\leftrow create\_for\_booking & \leftrow booking & \leftrow Классовый метод. Создаёт договор для бронирования или возвращает существующий. Обеспечивает наличие сгенерированного PDF-файла. \\ \hline
\leftrow generate\_document\_number & \leftrow --- & \leftrow Формирует уникальный номер документа на основе текущего года, идентификатора бронирования и фрагмента UUID. \\ \hline
\leftrow build\_preview\_text & \leftrow --- & \leftrow Формирует текстовый предпросмотр договора для отображения в интерфейсе. Содержит стороны, предмет, период, суммы и блоки ЭЦП в виде обычного текста. \\ \hline
\leftrow ensure\_generated\_file & \leftrow --- & \leftrow Вызывает \texttt{build\_pdf\_bytes}, сохраняет результат в файловое хранилище и фиксирует путь к файлу в базе данных. \\ \hline
\leftrow can\_be\_opened & \leftrow --- & \leftrow Проверяет, что статус бронирования допускает работу с договором (\texttt{confirmed} или \texttt{completed}). \\ \hline
\leftrow sign\_for\_user & \leftrow user, signer\_name, certificate\_pin & \leftrow Фиксирует подпись текущего пользователя: проверяет его роль, валидирует имя и PIN, записывает код и дату подписания. Если подписали оба участника, переводит договор в статус подписанного и перегенерирует PDF. \\ \hline
\leftrow generate\_signature\_code & \leftrow role & \leftrow Генерирует уникальный демонстрационный код электронной подписи в формате \texttt{EDS-DEMO-\{ROLE\}-\{booking\_id\}-\{UUID\}}. \\ \hline
\leftrow build\_pdf\_bytes & \leftrow --- & \leftrow Формирует полный PDF-файл договора с помощью ReportLab. Включает заголовок, преамбулу сторон, разделы договора, таблицу реквизитов и блоки ЭЦП. Возвращает байтовое содержимое файла. \\ \hline
\leftrow \_amount\_words & \leftrow amount & \leftrow Вспомогательный статический метод. Форматирует числовую сумму для вставки в текст договора. \\ \hline
\end{xltabular}

Описание представлений REST API приложения \texttt{contracts} представлено в таблице~\ref{tab:contracts_views}.

\begin{xltabular}{\textwidth}{|>{\small\raggedright\arraybackslash}p{0.30\textwidth}|>{\small\raggedright\arraybackslash}p{0.22\textwidth}|>{\small\raggedright\arraybackslash}X|}
\caption{Описание представлений REST API приложения contracts \label{tab:contracts_views}}\\ \hline
\centrow Метод & \centrow HTTP-метод, маршрут & \centrow Описание \\ \hline
\endfirsthead
\continuecaption{Продолжение таблицы \ref{tab:contracts_views}}
\centrow Метод & \centrow HTTP-метод, маршрут & \centrow Описание \\ \hline
\finishhead
\leftrow ContractViewSet.get\_queryset & \leftrow --- & \leftrow Формирует базовый QuerySet договоров, фильтруя по участию текущего пользователя в качестве арендатора или арендодателя. \\ \hline
\leftrow ContractViewSet.by\_booking & \leftrow GET /contracts/by-booking/\{booking\_id\}/ & \leftrow Возвращает или автоматически создаёт договор по идентификатору подтверждённого бронирования. Ответ содержит текстовый предпросмотр, ссылку на PDF, статус подписания и роль текущего пользователя. \\ \hline
\leftrow ContractViewSet.sign & \leftrow POST /contracts/\{id\}/sign/ & \leftrow Принимает имя подписанта и PIN-код, фиксирует подпись текущего пользователя. После подписания обеими сторонами перегенерирует PDF. Возвращает обновлённые данные договора. \\ \hline
\end{xltabular}

Описание сериализаторов приложения \texttt{contracts} представлено в таблице~\ref{tab:contracts_serializers}.

\begin{xltabular}{\textwidth}{|>{\small\raggedright\arraybackslash}p{0.30\textwidth}|>{\small\raggedright\arraybackslash}p{0.22\textwidth}|>{\small\raggedright\arraybackslash}X|}
\caption{Описание сериализаторов приложения contracts \label{tab:contracts_serializers}}\\ \hline
\centrow Сериализатор / поле & \centrow Тип & \centrow Описание \\ \hline
\endfirsthead
\continuecaption{Продолжение таблицы \ref{tab:contracts_serializers}}
\centrow Сериализатор / поле & \centrow Тип & \centrow Описание \\ \hline
\finishhead
\leftrow ContractSerializer & \leftrow ModelSerializer & \leftrow Полный сериализатор договора для чтения. Все поля доступны только на чтение. \\ \hline
\leftrow \quad file\_url & \leftrow SerializerMethodField & \leftrow Абсолютный URL PDF-файла договора с учётом домена запроса. \\ \hline
\leftrow \quad preview\_text & \leftrow SerializerMethodField & \leftrow Текстовый предпросмотр договора, формируемый методом \texttt{build\_preview\_text}. \\ \hline
\leftrow \quad status & \leftrow SerializerMethodField & \leftrow Статус договора: \texttt{draft}, \texttt{partially\_signed} или \texttt{signed}. \\ \hline
\leftrow \quad my\_role & \leftrow SerializerMethodField & \leftrow Роль текущего пользователя в договоре: \texttt{renter}, \texttt{owner} или \texttt{null}. \\ \hline
\leftrow \quad my\_signed & \leftrow SerializerMethodField & \leftrow Признак того, подписал ли текущий пользователь договор со своей стороны. \\ \hline
\leftrow ContractSignSerializer & \leftrow Serializer & \leftrow Сериализатор для приёма данных подписания договора. Поля: \texttt{signer\_name} и \texttt{certificate\_pin}. \\ \hline
\end{xltabular}

\subsection{Эксплуатация и развёртывание}

\subsubsection{Локальное развёртывание}
Проект рассчитан на запуск через Docker Compose. В составе инфраструктуры используются контейнеры приложения, базы данных PostgreSQL, Redis и reverse proxy. Такой подход упрощает развёртывание, так как окружение описано декларативно и не требует ручной настройки каждого сервиса на машине разработчика.

Базовый порядок запуска:
\begin{enumerate}
    \item Подготовить переменные окружения.
    \item Собрать контейнеры проекта.
    \item Запустить Docker Compose.
    \item Выполнить миграции базы данных.
    \item Создать пользователя администратора при необходимости.
    \item Проверить доступность API и WebSocket-маршрутов.
\end{enumerate}

\subsubsection{Параметры конфигурации}
Ключевые параметры окружения:
\begin{itemize}
    \item \texttt{SECRET\_KEY} --- секретный ключ Django;
    \item параметры подключения к PostgreSQL;
    \item параметры Redis;
    \item \texttt{EMAIL\_HOST} --- SMTP-сервер;
    \item \texttt{EMAIL\_PORT} --- порт SMTP;
    \item \texttt{EMAIL\_HOST\_USER} --- учётная запись отправителя;
    \item \texttt{EMAIL\_HOST\_PASSWORD} --- пароль приложения или токен;
    \item \texttt{MEDIA\_ROOT} --- директория хранения файлов, включая договоры и изображения сообщений;
    \item \texttt{MEDIA\_URL} --- URL-префикс для доступа к медиафайлам.
\end{itemize}

\subsubsection{Конфигурация технических средств}
Минимальная конфигурация для запуска проекта:
\begin{itemize}
    \item процессор с частотой от 2 ГГц;
    \item оперативная память от 4 ГБ;
    \item установленный Docker;
    \item доступ к сети для отправки email и работы внешних интеграций;
    \item свободное место на диске для базы данных и PDF-файлов.
\end{itemize}

\subsubsection{Наблюдаемость и сопровождение}
Для сопровождения разработанных модулей рекомендуется:
\begin{itemize}
    \item вести журнал ошибок отправки email;
    \item логировать ошибки WebSocket-подключений;
    \item контролировать количество непрочитанных уведомлений;
    \item отслеживать объём таблицы сообщений;
    \item периодически очищать устаревшие временные данные;
    \item выполнять резервное копирование PDF-договоров;
    \item проверять корректность фоновой команды напоминаний о возврате.
\end{itemize}

\subsection{Верификация технического проекта}

\subsubsection{Функциональная верификация}
Функциональная верификация должна подтверждать корректность пользовательских сценариев. Для подсистемы уведомлений проверяются следующие случаи:
\begin{itemize}
    \item при создании бронирования создаётся уведомление арендодателю;
    \item при подтверждении бронирования создаётся уведомление арендатору;
    \item при отмене бронирования уведомления получают обе стороны;
    \item при подтверждении оплаты создаётся уведомление арендодателю;
    \item при приближении даты возврата создаётся напоминание арендатору;
    \item пользователь может отметить одно уведомление как прочитанное;
    \item пользователь может отметить все уведомления как прочитанные;
    \item счётчик непрочитанных уведомлений возвращает корректное значение.
\end{itemize}

Для подсистемы чатов проверяются следующие случаи:
\begin{itemize}
    \item пользователь может создать диалог с владельцем товара;
    \item нельзя создать диалог с самим собой;
    \item диалоги по разным товарам разделяются;
    \item сообщения сохраняются в базе данных;
    \item пустое сообщение отклоняется;
    \item сообщение с изображением принимается;
    \item список диалогов сортируется по дате последнего сообщения;
    \item входящие сообщения отмечаются как прочитанные.
\end{itemize}

Для подсистемы договоров проверяются следующие случаи:
\begin{itemize}
    \item участник сделки может получить договор по подтверждённому бронированию;
    \item договор недоступен для бронирования в статусе ожидания;
    \item договор создаётся один раз для одного бронирования;
    \item арендатор может подписать договор;
    \item арендодатель может подписать договор;
    \item после подписания обеими сторонами договор получает статус подписанного;
    \item PDF-файл договора создаётся и доступен по ссылке.
\end{itemize}

\subsubsection{Техническая верификация}
Техническая верификация должна включать проверку API, WebSocket и операций с файлами.

Для API проверяются:
\begin{itemize}
    \item HTTP-коды успешных и ошибочных ответов;
    \item формат JSON-ответов;
    \item корректность сериализуемых полей;
    \item запрет доступа неавторизованным пользователям;
    \item запрет доступа к чужим объектам.
\end{itemize}

Для WebSocket проверяются:
\begin{itemize}
    \item отказ в подключении анонимному пользователю;
    \item отказ в подключении к чужому чату;
    \item доставка нового сообщения участникам диалога;
    \item доставка события прочтения сообщений;
    \item доставка нового уведомления в персональный канал пользователя.
\end{itemize}

Для PDF-договоров проверяются:
\begin{itemize}
    \item генерация файла при создании договора;
    \item пересоздание файла после подписания;
    \item наличие номера договора;
    \item наличие данных бронирования;
    \item наличие сведений о подписях сторон.
\end{itemize}

\subsection{Модульное и интеграционное тестирование разработанных модулей}
Модульное и интеграционное тестирование разработанных приложений выполнялось с использованием встроенных средств тестирования Django и Django REST Framework. Для проверки REST API применялся класс \texttt{APIClient}, позволяющий выполнять HTTP-запросы к API без запуска пользовательского интерфейса. Такой подход соответствует структуре разработанного проекта, так как реализованная часть представляет собой серверные Django-приложения и набор API.

Тестирование выполнялось в изолированной конфигурации. Для тестовой базы данных использовалась SQLite, WebSocket-каналы заменялись внутрипроцессным слоем \texttt{InMemoryChannelLayer}, а отправка email выполнялась через тестовый backend Django. Это позволило проверить бизнес-логику модулей без развёртывания внешней базы данных, Redis и SMTP-сервера.

Состав тестируемых модулей представлен в таблице~\ref{tab:module_tests}.

\begin{xltabular}{\textwidth}{|>{\small\raggedright\arraybackslash}p{0.22\textwidth}|>{\small\raggedright\arraybackslash}p{0.24\textwidth}|>{\small\raggedright\arraybackslash}X|}
\caption{Состав модульных и интеграционных тестов \label{tab:module_tests}}\\ \hline
\centrow Модуль & \centrow Файл тестов & \centrow Проверяемые функции \\ \hline
\endfirsthead
\continuecaption{Продолжение таблицы \ref{tab:module_tests}}
\centrow Модуль & \centrow Файл тестов & \centrow Проверяемые функции \\ \hline
\finishhead
\leftrow Чаты & \leftrow \texttt{chats/tests.py} & \leftrow Создание сообщений, проверка участников чата, отправка изображения, отклонение пустого сообщения, получение справки, отметка сообщений как прочитанных, разделение диалогов по товарам, сортировка диалогов по последнему сообщению. \\ \hline
\leftrow Уведомления & \leftrow \texttt{notifications/tests.py} & \leftrow Создание уведомления при бронировании, подтверждении бронирования, оплате, приближении даты возврата и новом сообщении; проверка служебных метаданных; проверка endpoint для отметки уведомления прочитанным; проверка вызова email-уведомлений. \\ \hline
\leftrow Договоры и электронная подпись & \leftrow \texttt{contracts/tests.py} & \leftrow Получение договора по подтверждённому бронированию, запрет создания договора для неподтверждённого бронирования, подписание договора арендатором и арендодателем с формированием демонстрационных кодов электронной подписи. \\ \hline
\end{xltabular}

Для запуска тестов использовалась команда:

\begin{verbatim}
python manage.py test chats notifications contracts --settings=arenda_site.test_settings --verbosity=2
\end{verbatim}

Результат запуска тестов представлен на рисунке~\ref{fig:module_tests_result}.

\begin{figure}[H]
\begin{lstlisting}[language=Python]
Found 19 test(s).
Creating test database for alias 'default' ('file:memorydb_default?mode=memory&cache=shared')...

test_chat_members (chats.tests.ChatTest.test_chat_members) ... ok
test_conversations_are_sorted_by_last_message (chats.tests.ChatTest.test_conversations_are_sorted_by_last_message) ... ok
test_create_conversation_is_separated_by_item (chats.tests.ChatTest.test_create_conversation_is_separated_by_item) ... ok
test_create_empty_message_via_api_is_rejected (chats.tests.ChatTest.test_create_empty_message_via_api_is_rejected) ... ok
test_create_message (chats.tests.ChatTest.test_create_message) ... ok
test_create_message_with_image_via_api (chats.tests.ChatTest.test_create_message_with_image_via_api) ... ok
test_mark_read_returns_updated_message_ids (chats.tests.ChatTest.test_mark_read_returns_updated_message_ids) ... ok
test_support_faq_is_available_without_auth (chats.tests.ChatTest.test_support_faq_is_available_without_auth) ... ok
test_booking_creation_creates_owner_notification (notifications.tests.NotificationTests.test_booking_creation_creates_owner_notification) ... ok
test_confirm_booking_creates_renter_notification (notifications.tests.NotificationTests.test_confirm_booking_creates_renter_notification) ... ok
test_legacy_return_reminder_command_creates_correct_notification (notifications.tests.NotificationTests.test_legacy_return_reminder_command_creates_correct_notification) ... Отправлено 1 напоминаний о возврате
ok
test_mark_notification_read_endpoints (notifications.tests.NotificationTests.test_mark_notification_read_endpoints) ... ok
test_new_message_notification_contains_chat_metadata (notifications.tests.NotificationTests.test_new_message_notification_contains_chat_metadata) ... ok
test_payment_confirmed_notification_points_to_my_items (notifications.tests.NotificationTests.test_payment_confirmed_notification_points_to_my_items) ... ok
test_payment_confirmed_sends_email (notifications.tests.NotificationTests.test_payment_confirmed_sends_email) ... ok
test_return_reminder_creates_notification (notifications.tests.NotificationTests.test_return_reminder_creates_notification) ... ok
test_contract_is_unavailable_for_pending_booking (contracts.tests.ContractApiTest.test_contract_is_unavailable_for_pending_booking) ... ok
test_participant_can_open_contract_by_booking (contracts.tests.ContractApiTest.test_participant_can_open_contract_by_booking) ... ok
test_renter_and_owner_can_sign_contract_with_demo_eds (contracts.tests.ContractApiTest.test_renter_and_owner_can_sign_contract_with_demo_eds) ... ok
----------------------------------------------------------------------
Ran 19 tests in 36.344s

OK
Destroying test database for alias 'default' ('file:memorydb_default?mode=memory&cache=shared')...	
\end{lstlisting}
\caption{Результат модульного и интеграционного тестирования приложений chats, notifications и contracts}
\label{fig:module_tests_result}
\end{figure}

Результаты тестирования приложения \texttt{chats} подтверждают корректность работы подсистемы обмена сообщениями. Были проверены создание сообщения, наличие участников диалога, отправка сообщения с изображением, отклонение пустого сообщения, получение справочной информации, отметка входящих сообщений как прочитанных, создание отдельных диалогов для разных товаров и сортировка списка диалогов по дате последнего сообщения.

Результаты тестирования приложения \texttt{notifications} подтверждают корректность создания уведомлений при ключевых событиях системы. Проверялись уведомления о новой заявке на бронирование, подтверждении бронирования, подтверждении оплаты, приближении даты возврата и получении нового сообщения. Также проверялась корректность метаданных уведомлений, так как они используются клиентской частью для перехода к нужному разделу системы.

Результаты тестирования приложения \texttt{contracts} подтверждают корректность работы подсистемы договоров. Были проверены получение договора по подтверждённому бронированию, запрет доступа к договору для бронирования в статусе ожидания, а также последовательное подписание договора арендатором и арендодателем. После подписания обеими сторонами договор получает статус подписанного, фиксируется дата подписания, формируются демонстрационные коды электронной подписи и пересоздаётся PDF-файл договора.

\subsection{Системное тестирование модуля взаимодействия с пользователем}

Системное тестирование проводилось на уровне серверного API и проверяло полный пользовательский сценарий взаимодействия арендатора и арендодателя. Поскольку в рамках работы разрабатывались Django-приложения и API, проверка выполнялась без пользовательского интерфейса, через последовательность операций серверной части.

Проверяемый системный сценарий включал следующие действия:

\begin{enumerate}
    \item Арендатор создаёт заявку на бронирование товара.
    \item Система создаёт уведомление арендодателю о новой заявке.
    \item Арендодатель подтверждает бронирование.
    \item Система создаёт уведомление арендатору о подтверждении бронирования.
    \item Арендатор создаёт или открывает чат с арендодателем.
    \item Арендатор отправляет сообщение в чат.
    \item Система сохраняет сообщение и создаёт уведомление получателю.
    \item По подтверждённому бронированию создаётся договор аренды.
    \item Арендатор подписывает договор демонстрационной электронной подписью.
    \item Арендодатель подписывает договор демонстрационной электронной подписью.
    \item После подписания обеими сторонами договор получает статус подписанного, а PDF-файл договора пересоздаётся.
\end{enumerate}

Результат системного тестирования представлен на рисунке~\ref{fig:system_tests_result}.

\begin{figure}[H]
\begin{lstlisting}[language=Python]
	Создание бронирования
	Уведомление арендодателю создано
	Подтверждение бронирования арендодателем
	Уведомление арендатору создано
	Создание диалога между участниками сделки
	Отправка сообщения в чат
	Уведомление о новом сообщении создано
	Cоздание договора по бронированию
	Подписание договора арендатором
	Подписание договора арендодателем
	Статус договора: подписан
	PDF-файл договора сформирован
\end{lstlisting}
\caption{Результат системного тестирования модуля взаимодействия с пользователем}
\label{fig:system_tests_result}
\end{figure}

В результате системного тестирования подтверждено, что разработанные приложения \texttt{chats}, \texttt{notifications} и \texttt{contracts} работают совместно и обеспечивают полный сценарий взаимодействия пользователей: от создания события сделки до коммуникации в чате, получения уведомлений, формирования договора и его подписания электронной подписью.